% --- //////////////////////////////////////////////////////////////// ---
% MODULE 2
\begin{longtable}{@{}p{4cm} p{11cm}@{}}
\caption{Mô tả Use Case: Chia sẻ Lộ trình}
\label{usecase-2-01} \\
\toprule
\textbf{Use case ID}       & UC-M2-01 \\ \midrule
\textbf{Tên Use case}      & Chia sẻ Lộ trình \\ \midrule
\textbf{Các tác nhân}      & Traveler \\ \midrule
\textbf{Mô tả}             & Cho phép Traveler chia sẻ một lộ trình đã lưu của mình cho một hoặc nhiều Traveler khác. Những người được chia sẻ chỉ có quyền xem nội dung lộ trình. \\ \midrule
\textbf{Tiền điều kiện}    & 
\begin{tabular}[t]{@{}p{10.5cm}@{}}
1. Traveler đã đăng nhập vào hệ thống.\\
2. Lộ trình muốn chia sẻ đã được lưu vào tài khoản.
\end{tabular} \\ \midrule
\textbf{Hậu điều kiện}     & Quyền xem lộ trình được cấp thành công cho các Traveler được chỉ định. \\ \midrule
\textbf{Luồng sự kiện chính} & 
\begin{tabular}[t]{@{}p{10.5cm}@{}}
1. Traveler mở một lộ trình đã lưu và nhấn vào nút ""Chia sẻ"".\\
2. Hệ thống hiển thị một giao diện, bao gồm một ô nhập liệu để tìm kiếm người dùng.\\
3. Traveler nhập tên tài khoản của người mà họ muốn chia sẻ. Hệ thống có thể hiển thị các gợi ý tìm kiếm.\\
4. Traveler chọn một hoặc nhiều tài khoản từ kết quả tìm kiếm.\\
5. Traveler nhấn nút ""Xác nhận chia sẻ"".\\
6. Hệ thống kiểm tra sự tồn tại và hợp lệ của các tài khoản được chọn.\\
7. Hệ thống tạo một bản ghi về quyền truy cập trong cơ sở dữ liệu, cấp quyền ""chỉ xem"" cho các tài khoản được chọn đối với lộ trình này.\\
8. Hệ thống hiển thị thông báo ""Đã chia sẻ lộ trình thành công!""
\end{tabular} \\ \midrule
\textbf{Luồng rẽ nhánh}    & N/A \\ \midrule
\textbf{Luồng ngoại lệ} &
\begin{tabular}[t]{@{}p{10.5cm}@{}}
\textbf{E1: Không tìm thấy tài khoản người dùng:}\\
Thời điểm: Xảy ra tại bước 3.\\
Mô tả: Traveler nhập một tên tài khoản không tồn tại trong hệ thống. Hệ thống sẽ hiển thị thông báo ""Không tìm thấy người dùng này.""\\[3pt]
\textbf{E2: Lỗi hệ thống khi cấp quyền}\\
Thời điểm: Xảy ra tại bước 7.\\
Mô tả: Hệ thống không thể lưu lại quyền chia sẻ do lỗi kết nối hoặc lỗi máy chủ. Hệ thống sẽ hiển thị thông báo ""Chia sẻ thất bại, vui lòng thử lại.""
\end{tabular} \\ 
\bottomrule
\end{longtable}

\begin{longtable}{@{}p{4cm} p{11cm}@{}}
\caption{Mô tả Use Case: Xem Lộ trình được chia sẻ}
\label{usecase-2-02} \\
\toprule
\textbf{Use case ID}       & UC-M2-02 \\ \midrule
\textbf{Tên Use case}      & Xem Lộ trình được chia sẻ \\ \midrule
\textbf{Các tác nhân}      & Traveler \\ \midrule
\textbf{Mô tả}             & Cho phép Traveler xem danh sách các lộ trình mà những người dùng khác đã chia sẻ cho họ, và mở ra để xem chi tiết. \\ \midrule
\textbf{Tiền điều kiện}    & 
\begin{tabular}[t]{@{}p{10.5cm}@{}}
1. Traveler đã đăng nhập vào hệ thống.\\
2. Có ít nhất một Traveler khác đã chia sẻ lộ trình cho họ.
\end{tabular} \\ \midrule
\textbf{Hậu điều kiện}     & Traveler xem được nội dung chi tiết của lộ trình đã được chia sẻ. \\ \midrule
\textbf{Luồng sự kiện chính} &
\begin{tabular}[t]{@{}p{10.5cm}@{}}
1. Traveler truy cập vào trang quản lý hành trình cá nhân.\\
2. Họ chọn vào mục hoặc tab ""Được chia sẻ với tôi"".\\
3. Hệ thống truy vấn và hiển thị một danh sách các lộ trình mà Traveler này đã được cấp quyền xem, có thể kèm theo tên người chia sẻ.\\
4. Traveler nhấn vào một lộ trình trong danh sách.\\
5. Hệ thống hiển thị nội dung chi tiết của lộ trình đó ở chế độ chỉ xem (view-only). Tất cả các nút chỉnh sửa, xóa đều bị ẩn hoặc vô hiệu hóa.
\end{tabular} \\ \midrule
\textbf{Luồng rẽ nhánh}    & N/A \\ \midrule
\textbf{Luồng ngoại lệ} &
\begin{tabular}[t]{@{}p{10.5cm}@{}}
\textbf{E1: Người chia sẻ đã thu hồi quyền truy cập:}\\
Thời điểm: Xảy ra tại bước 4.\\
Mô tả: Traveler nhấn vào một lộ trình nhưng người chia sẻ đã hủy quyền truy cập trước đó. Hệ thống sẽ hiển thị thông báo ""Bạn không còn quyền truy cập vào lộ trình này.""\\[3pt]
\textbf{E2: Lộ trình gốc đã bị xóa:}\\
Thời điểm: Xảy ra tại bước 4.\\
Mô tả: Lộ trình mà người dùng muốn xem đã bị người tạo ra xóa vĩnh viễn. Hệ thống sẽ hiển thị thông báo ""Lộ trình này không còn tồn tại.""
\end{tabular} \\
\bottomrule
\end{longtable}

\begin{longtable}{@{}p{4cm} p{11cm}@{}}
\caption{Mô tả Use Case: Hủy Lộ trình}
\label{usecase-2-04} \\
\toprule
\textbf{Use case ID}       & UC-M2-04 \\ \midrule
\textbf{Tên Use case}      & Hủy Lộ trình \\ \midrule
\textbf{Các tác nhân}      & Traveler \\ \midrule
\textbf{Mô tả}             & Cho phép Traveler thay đổi trạng thái của một lộ trình từ ""Đang tiến hành"" sang ""Đã hủy"". Traveler sẽ được yêu cầu cung cấp lý do cho việc hủy. \\ \midrule
\textbf{Tiền điều kiện}    & 
\begin{tabular}[t]{@{}p{10.5cm}@{}}
1. Traveler đã đăng nhập.\\
2. Traveler có một lộ trình đang ở trạng thái ""Đang tiến hành""
\end{tabular} \\ \midrule
\textbf{Hậu điều kiện}     & Trạng thái của lộ trình được chuyển thành ""Đã hủy"". Lý do hủy (nếu có) được ghi nhận. \\ \midrule
\textbf{Luồng sự kiện chính} &
\begin{tabular}[t]{@{}p{10.5cm}@{}}
1. Vào ngày bắt đầu của lộ trình, hệ thống tự động chuyển trạng thái thành ""Đang tiến hành"".\\
2. Giao diện người dùng hiển thị lộ trình này kèm theo nút ""Hủy Lộ trình"".\\
3. Traveler nhấn vào nút ""Hủy Lộ trình"".\\
4. Hệ thống hiển thị một hộp thoại, yêu cầu người dùng nhập lý do hủy (có thể là một ô văn bản hoặc chọn từ danh sách có sẵn).\\
5. Traveler nhập lý do và nhấn ""Xác nhận hủy"".\\
6. Hệ thống cập nhật trạng thái của lộ trình trong cơ sở dữ liệu thành ""Đã hủy"".\\
7. Hệ thống lưu lại lý do mà người dùng đã cung cấp (phục vụ cho mục đích thống kê sau này).\\
8. Giao diện hiển thị thông báo ""Lộ trình đã được hủy"".
\end{tabular} \\ \midrule
\textbf{Luồng rẽ nhánh}    &
\begin{tabular}[t]{@{}p{10.5cm}@{}}
\textbf{A1: Hủy không cần lý do:}\\
Thời điểm: Xảy ra tại bước 5, nếu việc nhập lý do là không bắt buộc.\\
Mô tả: Traveler có thể bỏ qua việc nhập lý do và nhấn ""Xác nhận hủy"" trực tiếp. Luồng tiếp tục bình thường.
\end{tabular} \\ \midrule
\textbf{Luồng ngoại lệ} &
\begin{tabular}[t]{@{}p{10.5cm}@{}}
\textbf{E1: Lỗi kết nối hoặc lỗi máy chủ:}\\
Thời điểm: Xảy ra tại bước 6.\\
Mô tả: Hệ thống không thể cập nhật trạng thái do lỗi. Hệ thống sẽ hiển thị thông báo ""Hủy thất bại, vui lòng thử lại.""
\end{tabular} \\
\bottomrule
\end{longtable}

\begin{longtable}{@{}p{4cm} p{11cm}@{}}
\caption{Mô tả Use Case: Tích hợp Lộ trình với Google Calendar}
\label{usecase-2-06} \\
\toprule
\textbf{Use case ID}       & UC-M2-06 \\ \midrule
\textbf{Tên Use case}      & Tích hợp Lộ trình với Google Calendar \\ \midrule
\textbf{Các tác nhân}      & Traveler (Primary), Google Calendar (Secondary) \\ \midrule
\textbf{Mô tả}             & Cho phép Traveler đồng bộ (xuất) một lộ trình đã lưu của họ sang tài khoản Google Calendar cá nhân để tiện theo dõi lịch trình trên các thiết bị khác. \\ \midrule
\textbf{Tiền điều kiện}    & 
\begin{tabular}[t]{@{}p{10.5cm}@{}}
1. Traveler đã đăng nhập vào hệ thống.\\
2. Traveler có ít nhất một lộ trình đã lưu (có thông tin ngày, giờ cụ thể).\\
3. Traveler đang xem chi tiết lộ trình mà họ muốn đồng bộ.
\end{tabular} \\ \midrule
\textbf{Hậu điều kiện}     & 
\begin{tabular}[t]{@{}p{10.5cm}@{}}
(Thành công): Các hoạt động trong lộ trình được tạo thành công dưới dạng các Sự kiện (Events) trên Google Calendar của Traveler.\\
(Thất bại): Lộ trình không được đồng bộ, hệ thống hiển thị thông báo lỗi.
\end{tabular} \\ \midrule
\textbf{Luồng sự kiện chính} &
\begin{tabular}[t]{@{}p{10.5cm}@{}}
1. Traveler đang xem chi tiết một lộ trình đã lưu.\\
2. Traveler nhấn vào nút “Tích hợp Google Calendar” (hoặc “Thêm vào Lịch”).\\
3. Hệ thống chuyển hướng Traveler đến trang xác thực OAuth 2.0 của Google.\\
4. Hệ thống yêu cầu Traveler cấp quyền cho ứng dụng TravelEZ để ""Xem và quản lý lịch"" của họ.\\
5. Traveler chọn tài khoản Google và nhấn ""Cho phép"".\\
6. Google trả về mã xác thực (authorization code) cho hệ thống TravelEZ.\\
7. Hệ thống dùng mã này để lấy Access Token truy cập Google Calendar API.\\
8. Hệ thống đọc dữ liệu chi tiết của lộ trình (tên hoạt động, địa điểm, thời gian bắt đầu, thời gian kết thúc).\\
9. Hệ thống lặp qua từng hoạt động và gọi API của Google Calendar để tạo các sự kiện tương ứng.\\
10. Hệ thống hiển thị thông báo ""Đã đồng bộ lộ trình sang Google Calendar thành công!""
\end{tabular} \\ \midrule
\textbf{Luồng rẽ nhánh}    &
\begin{tabular}[t]{@{}p{10.5cm}@{}}
\textbf{A1: Người dùng đã cấp quyền trước đó:}\\
Tại bước 3, hệ thống phát hiện đã có Access Token hợp lệ của người dùng.\\
Hệ thống bỏ qua các bước 4, 5, 6, 7.\\
Luồng sự kiện tiếp tục từ bước 8.
\end{tabular} \\ \midrule
\textbf{Luồng ngoại lệ} &
\begin{tabular}[t]{@{}p{10.5cm}@{}}
\textbf{E1: Người dùng từ chối cấp quyền:}\\
Tại bước 5, Traveler nhấn ""Từ chối"" hoặc đóng cửa sổ xác thực.\\
Hệ thống hiển thị thông báo ""Không thể tích hợp nếu không được cấp quyền truy cập."" Use case kết thúc.\\[3pt]
\textbf{E2: Lỗi Google Calendar API:}\\
Tại bước 9, API của Google trả về lỗi (ví dụ: token hết hạn, API service down, không tìm thấy lịch chính...).\\
Hệ thống hiển thị thông báo ""Đã xảy ra lỗi khi đồng bộ. Vui lòng thử lại sau.""\\[3pt]
\textbf{E3: Lộ trình không có dữ liệu thời gian:}\\
Tại bước 8, hệ thống phát hiện lộ trình đang xem không có thông tin ngày/giờ cụ thể (ví dụ: chỉ là một danh sách địa điểm).\\
Hệ thống hiển thị thông báo ""Không thể đồng bộ lộ trình chưa được lập lịch. Vui lòng thêm ngày giờ cho các hoạt động."" Use case kết thúc.
\end{tabular} \\
\bottomrule
\end{longtable}